% Transcription — fonds Alexandre Grothendieck, shelfmark 59, batch 1 (pages 1-20).
% Established from the facsimile archives/batches/59/batch-01.pdf (PDF page k =
% archivists' page k, no cover sheet), cut from a copy of 59.pdf that was
% fetched through the project's own relay
% (https://grothendieck-relay.onrender.com/source/59.pdf), not uploaded by the
% user; tiles in archives/tiles/59/p1-p20. Per the skill transcribe-grothendieck.
% The folder is twenty-seven pages; this is the first of two batches (pages
% 21-27 are batch 2, transcribed separately; its file was read before this one
% was finished and the Picard notation is the same in both).
%
% Pass: Opus 5.5 (claude-opus-5-5), 2026-09-26 — first pass, unchecked against the pages by a human.
%
% Pages skipped: none. Pages summarised: none. Pages transcribed: 1-20, twenty
% leaves in blue-black ink with a few pencil additions (pp. 1, 4, 6). His own
% numbers stand in the top right corner and, where legible (pp. 16-20 at
% least), coincide with the archivists' numbering, so no \note records them.
% This run is distinct from the run paginated 1-6 on pp. 22-27 (batch 2).
% Fast drafting throughout: the formulas carry, much of the connective prose
% and most oblique margin notes do not and are left \ill{}. Pages 11, 12, 18
% and 20 carry framed blocks cancelled by long diagonals or loops; they are
% given as \struck{} with the few words that can be read.
%
% What the batch is about. X/S proper, flat, of finite presentation,
% f_*(O_X) = O_S universally, P = Pic_{X/S}, P^0 = B an abelian scheme,
% A = Pic^0_{B/S} (= Alb^0_{X/S}). (a) For a section g, the canonical sheaf
% L_g on X x_S P and M_g = det^* R f_{P*}(L_g); (b) change of section,
% M_{g'} = M_g (N_g N_{g'}^{-1})^chi, chi the Euler characteristic; (c) hence
% phi : NS_{X/S} -> NS_{B/S}, independent of g, by descent without a section;
% (d) Pic_{P^delta/S} as the twist of Pic_{B/S} by the torsor P^delta, and a
% canonical section L(delta) of the A-torsor Alb^{chi(delta)} x^A
% Pic^{phi(delta)}_{B/S} x^A (P^delta x^B (A, phi(delta)~)), with u_delta =
% -chi(delta) id_A (« vérifier si le signe est correct ! »). (e) Computation
% of phi(delta) by Riemann-Roch / Chern character, l-adically: phi(delta) =
% (1/2) pr_{2*}(pr_1^*(a_{n-1}) D^2); Ex. 1, curves: phi(delta) is the canonical
% polarisation of the Jacobian, chi = 1 - g + d, the torsor becomes P^{2g-2}
% and the expected section is K; N.B. on the moduli of curves. (f) Ex. 2,
% abelian schemes: phi(delta) = Phi(Delta)^{*(n-1)}/(n-1)! via the Fourier
% transform on Chow groups, phi(delta)~ delta~ = -chi(delta) id,
% chi(delta') = chi(delta)^{n-1}, phi(phi(delta)) = delta chi(delta)^{n-2};
% canonical isomorphism (*) Pic^{delta'}_{B/S} = Pic^delta_{A/S} x^B (A, delta'~);
% consequences: (Pic^{delta'}_{B/S})^{(x)2} and (Pic^delta)^{(x) 2 chi^{n-2}}
% canonically trivial, and more when chi(delta) is even (pp. 12-13).
% Pp. 14-16: the same construction twisted by a perfect complex F on X,
% phi_F additive in F, computed for curves (chi_F = 1 - g + d + c, phi_F the
% canonical polarisation; « pas de miracle ») and abelian schemes. Pp. 16-20,
% a second paragraph lettered f): the homomorphism l' : X x Pic_{X/S} ->
% Pic_{B/S} built from the Weil sheaf W(g), its dependence on g through
% psi_X : X -> Alb^1, the formula l_g(lambda + beta) = l_g(lambda) +
% phi(delta)~(beta), proved via det F_L(L (x) L_beta) (a Fourier-type
% transform, letter doubtful), and the identification with l_delta of d);
% a Remarque on replacing the section by a given hom. B -> Pic_{X/S} and a
% torsor Q under A (« petite topologie de Albanese », reading doubtful).
%
% Notation fixed once for the batch: \underline{\mathrm{Pic}},
% \underline{\mathrm{Alb}}, \underline{\mathrm{NS}} as in batch 2 and the rest
% of the fonds (he underlines all three); P^\delta, P^0 = B; \mathcal{L}_g for
% the sheaf of a) (the page also writes L_g on p. 3, kept); the contracted
% products X \overset{A}{\times} Y as he stacks the group over the cross;
% \widetilde{\varphi(\delta)} for the map B -> A defined by phi(delta);
% \mathfrak{F}_{\mathcal{L}} for the curly letter of pp. 18-19 (doubtful).
% He letters two paragraphs e) (pp. 5 and 14) and two f) (pp. 9 and 16); both
% are kept, the first f) being « Ex. 2 ».
%
% Readings to check first: the sign and exponent of Alb in the boxed formula
% of p. 4 (chi(delta) vs -chi(delta)); the exponents in lambda on p. 12; the
% « p.p.c.m. » passage of p. 13; the curly letter F_L on pp. 18-19; the arrow
% label tau and the bracket « petite topologie de Albanese » on p. 20; all
% oblique margin notes (pp. 1, 10, 11, 12, 18, 19), which are mostly \ill{}.

\documentclass[11pt,a4paper]{article}
% Le préambule commun à toutes les transcriptions du fonds Grothendieck.
%
% Il tient en une idée : **l'incertitude se compose, elle ne se résout pas.**
% Une transcription qui lisse un mot douteux a détruit la seule chose pour
% laquelle on la faisait — un lecteur doit pouvoir savoir, à chaque endroit,
% ce qui était sur le feuillet et ce qui a été ajouté. Les six macros qui
% suivent sont donc l'appareil critique, et non de la décoration.
%
% Les mêmes commandes sont reconnues par `scripts/render.mjs`, qui produit la
% vue de lecture du site. Une macro ajoutée ici doit l'être aussi là-bas,
% faute de quoi le rendu échouera bruyamment — ce qui est voulu : un rendu
% silencieusement faux est pire qu'un rendu absent.

\ProvidesPackage{grothendieck}[2026/08/08 Transcriptions du fonds Grothendieck]

\usepackage{amsmath,amssymb,amsthm}
\usepackage{mathrsfs}
\usepackage{stmaryrd}
% Les diagrammes commutatifs sont partout dans ce fonds : c'est la langue
% même de la géométrie algébrique de Grothendieck, et une transcription qui
% les rendrait en prose ne transcrirait rien.
\usepackage{tikz-cd}
% Français par défaut — c'est la langue des feuillets — mais l'anglais est
% chargé aussi : la traduction et le résumé sont des documents anglais, et la
% typographie française (espaces avant les deux-points) y serait une faute.
% Ces documents basculent par \selectlanguage{english} en tête de corps.
\usepackage[english,french]{babel}
\usepackage{fontspec}
\usepackage{geometry}
\geometry{a4paper,margin=25mm,marginparwidth=28mm}
\usepackage{marginnote}
\usepackage{xcolor}
% ulem, not soul. `soul`'s \st and \ul are built on a character-by-character
% scan that XeLaTeX's font machinery breaks: the document fails with an
% undefined control sequence rather than degrading. ulem does the same two jobs
% and is Unicode-engine safe. `normalem` stops it hijacking \emph, which the
% transcriptions use for Grothendieck's own emphasis.
\usepackage[normalem]{ulem}
\usepackage{hyperref}
\hypersetup{colorlinks=true,linkcolor=black,urlcolor=black,pdfborder={0 0 0}}

% --- Métadonnées du lot -----------------------------------------------------
% Reprises telles quelles par le script de rendu, qui les affiche en tête de
% la vue de lecture. Elles fixent ce que le fichier prétend couvrir : sans
% elles, une transcription flotte sans point d'attache dans un fonds de seize
% mille feuillets.

\newcommand{\folder}[1]{\gdef\grothFolder{#1}}
\newcommand{\batch}[1]{\gdef\grothBatch{#1}}
\newcommand{\leaves}[2]{\gdef\grothFirst{#1}\gdef\grothLast{#2}}
\newcommand{\foldertitle}[1]{\gdef\grothFoldertitle{#1}}
\gdef\grothFolder{?}\gdef\grothBatch{?}\gdef\grothFirst{?}\gdef\grothLast{?}\gdef\grothFoldertitle{}

% --- Le résumé --------------------------------------------------------------
%
% Il ouvre la lecture modernisée, sous le titre « Résumé », et il est composé
% en petit. Ce n'est pas un ornement : c'est le seul passage du document qui
% ne soit pas des mathématiques, et un lecteur doit pouvoir le reconnaître
% avant de l'avoir lu — puis le sauter s'il connaît déjà le sujet, ou s'y
% arrêter s'il ne le connaît pas. Le filet qui le suit marque la reprise du
% corps.

\newenvironment{resume}
  {\small\setlength{\parskip}{0.45em}}
  {\par\medskip\hrule\medskip}

% --- Mots-clés ---------------------------------------------------------------
%
% En anglais, en clôture du résumé de la lecture modernisée : le vocabulaire
% moderne sous lequel chercher ce que les feuillets construisent. Le manifeste
% du site (`npm run manifest`) les extrait pour étiqueter la cote entière —
% cette ligne est la source unique des tags, il n'y a pas de fichier à côté.
\newcommand{\keywords}[1]{\par\smallskip\noindent
  {\footnotesize\textsc{Keywords} --- \textit{#1}}\par}

% --- L'appareil critique ----------------------------------------------------

% Le numéro de feuillet, en marge. C'est l'ancre qui permet de régler un
% désaccord sur une formule en regardant *un* feuillet plutôt qu'une chemise
% de six cents. Le site s'en sert aussi pour tourner les pages du fac-similé
% quand on fait défiler la transcription.
\newcommand{\leaf}[1]{%
  \marginnote{\footnotesize\color{gray!70}\textsf{#1}}%
  \label{leaf:#1}\ignorespaces}

% Illisible : on ne devine pas. Un mot inventé qui se lit comme les autres est
% le pire résultat possible.
\newcommand{\ill}{\textcolor{red!60!black}{[\,\ldots\,]}}

% Lecture douteuse : le mot est donné, et signalé comme incertain.
\newcommand{\uncertain}[1]{\uline{#1}}

% Ajout de l'éditeur — une lettre manquante, un mot sous-entendu.
\newcommand{\add}[1]{\textcolor{blue!50!black}{[#1]}}

% Biffé par Grothendieck lui-même. Ce qu'il a rayé fait partie du document :
% une rature dit souvent où la pensée a changé de direction.
\newcommand{\struck}[1]{\sout{#1}}

% Note du transcripteur, sur l'état matériel du feuillet.
\newcommand{\note}[1]{\par\smallskip{\small\color{gray!60!black}\textit{[#1]}}\par\smallskip}

% Note marginale de Grothendieck, distincte du corps du texte.
\newcommand{\marginal}[1]{\par\smallskip\noindent\hspace{1em}%
  {\small\color{gray!60!black}#1}\par\smallskip}

% --- Titre ------------------------------------------------------------------

\AtBeginDocument{%
  \begin{center}
    {\large\bfseries Fonds Alexandre Grothendieck --- cote n\textsuperscript{o}~\grothFolder}\\[2pt]
    {\itshape\grothFoldertitle}\\[4pt]
    {\small feuillets \grothFirst--\grothLast\ (lot \grothBatch)}\\[2pt]
    {\footnotesize\color{gray!60!black}Transcription établie d'après le fac-similé numérisé
     par l'Université de Montpellier.\\
     Les lectures incertaines sont soulignées, les illisibles notées [\,\ldots\,],
     les ajouts entre crochets.}
  \end{center}
  \vspace{1em}}

\endinput


\folder{59}
\batch{1}
\pages{1}{20}
\dating{[années 1960-1970]}
\watermark{Édition de démonstration}
\foldertitle{Construction des faisceaux inversibles sur schémas de Picard : notes manuscrites (s.d.)}

\begin{document}

\page{1}
\begin{tikzcd}
X \arrow[d, "f"'] \\
S
\end{tikzcd}

$f$ propre plat présentation finie, $f_{*}(\mathcal{O}_X) = \mathcal{O}_S$ universel !

\[ \underline{\mathrm{Pic}}_{X/S} = P \]

\begin{tikzcd}
X \arrow[r, no head] \arrow[d, no head] & X_P \arrow[d, no head] \\
S \arrow[r, no head] & P
\end{tikzcd}

\note{carré dessiné en haut à droite, sans pointes de flèches. La lettre $P$ est tracée tantôt ronde (une ronde), tantôt droite ; de même pour $B$ plus bas : on écrit $P$ et $B$ partout}

a) Supposons $X/S$ muni d'une section $g$. Alors sur $X \times_S P$ il y a
un faisceau inversible canonique $\mathcal{L}_g$ [trivialisé sur
$g(S) \times_S P$]. On en conclut
\[
(Rf_P)_{*}(\mathcal{L}_g) \in D_{\mathrm{parf}}(P)
\qquad [\, f_P : X \times_S P \to P \,].
\]
Nous nous intéressons en particulier :
\[
{\det}^{*}(Rf_P)_{*}(\mathcal{L}_g) = M_g
\]
qui est un module inversible [\uncertain{gradué} \ill{}] sur $P$.

\marginal{\ill{} en $S$ \ill{} [\uncertain{projectif}]}\note{note oblique dans la marge gauche, en face de la formule, presque entièrement illisible}

b) Remplaçons $g$ par une autre section $g'$. On a alors un isom
canonique
\[
\mathcal{L}_{g'} \simeq \mathcal{L}_g \otimes_{\mathcal{O}_P} g'^{*}_P(\mathcal{L}_g)^{-1}
\]
d'où :
\[
(Rf_P)_{*}(\mathcal{L}_{g'}) \simeq (Rf_P)_{*}(\mathcal{L}_g)\, g'^{*}_P(\mathcal{L}_g)^{-1}
\]
et par suite un isom canonique
\[
(*) \qquad M_{g'} \simeq M_g \, g'^{*}_P(\mathcal{L}_g)^{-\chi} = M_g \,(N_g N_{g'}^{-1})^{\chi}
\]
\note{au crayon, deux ovales entourent $g'^{*}_P(\mathcal{L}_g)^{-1}$ dans la ligne précédente et $g'^{*}_P(\mathcal{L}_g)$ ici ; un trait de crayon souligne $N_g N_{g'}$}

\marginal{$N_g N_{g'}^{-1}$ est algébriquement \emph{équivalent} à zéro sur $P$ rel$/S$ ; $M_{g'}$ et $M_g$ alg.\ équiv.\ : $M_g$ \uncertain{polaris\ldots}}\note{écrit verticalement dans la marge gauche, encadré, une flèche le rattache au début de c)}

où $\chi$ est la fonction \uncertain{localement constante} $P \to \mathbb{Z}$
donnée par le rang de $Rf_P(\mathcal{L}_g)$, i.e.\ la car.\ d'EP des faisceaux
inversibles \struck{sur W} [sur] les fibres de $X$ \ldots\ldots, et où $N$
désigne un faisceau inv.\ \add{quelconque sur} $X \times_S P$, et $N_g$
(resp.\ $N_{g'}$) le faisceau inv.\ sur $P$ qu'il définit via la section $g$
(resp.\ $g'$). On notera que le facteur de droite \ill{}.
\note{« quelconque sur » est inséré dans l'interligne}

c) Soit $\delta$ une section de $\underline{\mathrm{NS}}_{X/S}$, qui définit
donc un torseur $P^{\delta}$ sous $P^{0} = \underline{\mathrm{Pic}}^{0}_{X/S}$,
et considérons la restriction de $M_g$ à $P^{\delta}$. Nous supposons, pour
simplifier, par la suite, que $P^{0}$ est un schéma \emph{abélien} sur $S$.
\note{au-dessus de « $M_g$ », une insertion « $M_g$ \ldots » que le trait rattache à $P^{\delta}$}

\page{2}

Notons que $M_g$ et $M_{g'}$ définissent la même section de
$\underline{\mathrm{NS}}_{P^{\delta}/S}(S)$ en vertu de b), \uncertain{or} on a
un isom canonique
\[
\underline{\mathrm{NS}}_{P^{\delta}/S} \simeq \underline{\mathrm{NS}}_{P^{0}/S} .
\]
Donc on trouve une section
\[
\varphi(\delta) \in \underline{\mathrm{NS}}_{P^{0}/S}(S)
\]
indépendante du choix de la section $g$ de $X/S$. Par
\marginal{\emph{détailler}}
descente, elle est définie sans supposer qu'il existe une section. Comme la
formation de $\varphi(\delta)$ est compatible (par construction) à l'extension
de la base, on trouve en fait un homomorphisme de schémas (pas nécessairement
additif !)
\[
\boxed{\varphi : \underline{\mathrm{NS}}_{X/S} \longrightarrow \underline{\mathrm{NS}}_{B/S}}
\qquad (B = P^{0} = \underline{\mathrm{Pic}}^{0}_{X/S}).
\]

d) Revenons à $\delta$ \struck{\ill{}}. Notons que
$\underline{\mathrm{Pic}}_{P^{\delta}/S}$ se déduit de
$\underline{\mathrm{Pic}}_{B/S}$ en tordant par \struck{l'action} le torseur
$P^{\delta}$ sous $P^{0}$, via l'action naturelle de $B$ sur
$\underline{\mathrm{Pic}}_{B/S}$ déduite de l'action de $B$ sur lui-même par
translations. \struck{\ill{}} Cette action est triviale sur
$\underline{\mathrm{Pic}}^{0}_{B/S}$ et $\underline{\mathrm{NS}}_{B/S}$ [d'où
en particulier des isom
\[
\underline{\mathrm{Pic}}^{0}_{P^{\delta}/S} \simeq \underline{\mathrm{Pic}}^{0}_{B/S} ,
\qquad
\underline{\mathrm{NS}}_{P^{\delta}/S} \simeq \underline{\mathrm{NS}}_{B/S} \;].
\]

\page{3}

Si $\lambda$ est une section de $\underline{\mathrm{NS}}_{B/S}$ [donc définit
une section, notée encore $\lambda$, de
$\underline{\mathrm{NS}}_{P^{\delta}/S}$], alors
$\underline{\mathrm{Pic}}^{\lambda}_{P^{\delta}/S}$ se déduit de
$\underline{\mathrm{Pic}}^{\lambda}_{B/S}$ en tordant à l'aide de
$P^{\delta}$. Or soit
\[
A = \underline{\mathrm{Pic}}^{0}_{B/S} \quad (= \underline{\mathrm{Alb}}^{0}_{X/S})
\]
le schéma abélien dual de $B$, \struck{\ill{}} alors $\lambda$ définit un hom
\[
\tilde{\lambda} : B \longrightarrow A
\]
et l'action de $B$ sur $\underline{\mathrm{Pic}}^{\lambda}_{B/S}$ se déduit de
l'action évidente de $A = \underline{\mathrm{Pic}}^{0}_{B/S}$ sur
$\underline{\mathrm{Pic}}^{\lambda}_{B/S}$ (qui est un $A$-torseur) via l'hom
$\tilde{\lambda}$. Donc on trouve
\[
\underline{\mathrm{Pic}}^{\lambda}_{P^{\delta}/S} \simeq
\underline{\mathrm{Pic}}^{\lambda}_{B/S} \overset{A}{\times}
\underbrace{\bigl(P^{\delta} \overset{B}{\times} A\bigr)}
\]
\marginal{$A$-torseur déduit du $B$-torseur $P^{\delta}$ grâce à l'hom.\ $\tilde{\lambda} : B \to A$}\note{écrit sous l'accolade}

\struck{\ill{}} Reprenant les notations de c), nous appliquons ceci au cas où
l'on a
\[ \lambda = \varphi(\delta) , \]
on trouve que, pour toute \struck{\ill{}} section $g$ de $X/S$, la
\struck{restriction} $L_g \mid P^{\delta}$ définit une section de
$\underline{\mathrm{Pic}}^{\lambda}_{P^{\delta}/S}$, d'où une section
\note{la page porte bien $L_g$, là où c) parlait de $M_g$}

\page{4}

\[
\ell(g) \in \underline{\mathrm{Pic}}^{\varphi(\delta)}_{B/S} \overset{A}{\times}
\bigl(P^{\delta} \overset{B}{\times} (A, \tilde{\lambda})\bigr)
\]
La formation de cette dernière commute à tout changement de base, donc on
définit un morphisme de schémas (\uncertain{ni plat}, \ill{} \ldots)
\[
X \xrightarrow{\ \ell_{\delta}\ } \underline{\mathrm{Pic}}^{\varphi(\delta)}_{B/S}
\overset{A}{\times} \bigl(P^{\delta} \overset{B}{\times} (A, \widetilde{\varphi(\delta)})\bigr) ,
\]
\note{dans le second facteur, $\widetilde{\varphi(\delta)}$ est écrit par-dessus un premier symbole ; au crayon, en marge gauche, un petit schéma entouré : $A \leftarrow B$, avec $A = \underline{\mathrm{Alb}}^{0}$ et $B = \underline{\mathrm{Pic}}^{0}$ en dessous ; deux ovales au crayon entourent les deux membres de la flèche}
\uncertain{Comme} \ill{} est un torseur sous un schéma abélien $A$, on trouve
\add{une factorisation de $\ell_{\delta}$ par} un homomorphisme
\[
\underline{\mathrm{Alb}}^{1}_{X/S} \struck{(X/S)} \longrightarrow
\underline{\mathrm{Pic}}^{\varphi(\delta)}_{B/S} \overset{A}{\times}
\bigl(P^{\delta} \overset{B}{\times} (A, \widetilde{\varphi(\delta)})\bigr)
\]
qui est \struck{un homomorphisme} \add{compatible avec un} certain homomorphisme
\[
u_{\delta} : A \longrightarrow A
\]
sur les groupes structuraux,
$A = \underline{\mathrm{Alb}}^{0}_{X/S}$ \struck{\ill{}}. Pour voir quel il est,
on vérifie la formule
\[
u_{\delta} = -\chi(\delta)\, \mathrm{id}_A ,
\]
où $\chi = \chi(\delta)$ est \uncertain{(de b))} la fonction $S \to \mathbb{Z}$
dont l'image inverse sur $P^{\delta}$ est $\chi \mid P^{\delta}$ ($\chi$ comme
dans b)). On trouve donc un \struck{A} homomorphisme de $A$-torseurs
\note{l'interligne « une factorisation de $\ell_{\delta}$ par » et « compatible avec un » sont des insertions ; au-dessus de la seconde flèche, un « $A$ » au crayon}
\marginal{vérifier si le signe est correct !}
\[
\underline{\mathrm{Alb}}^{-\chi(\delta)}_{X/S} \longrightarrow
\underline{\mathrm{Pic}}^{\varphi(\delta)}_{B/S} \overset{A}{\times}
\bigl(P^{\delta} \overset{B}{\times} (A, \widetilde{\varphi(\delta)})\bigr)
\]
i.e.\ une \emph{section canonique}
\[
\boxed{L(\delta) \in \Gamma\Bigl(S, \underline{\mathrm{Alb}}^{\chi(\delta)}_{X/S}
\overset{A}{\times} \underline{\mathrm{Pic}}^{\varphi(\delta)}_{B/S}
\overset{A}{\times} \bigl(P^{\delta} \overset{B}{\times} (A, \widetilde{\varphi(\delta)})\bigr)\Bigr)}
\]
\note{dans le cadre, l'exposant de $\underline{\mathrm{Alb}}$ est $\chi(\delta)$, sans le signe moins de la ligne précédente ; on le laisse tel quel. Sous le second $\times$, un « $X/S$ » au crayon}

\page{5}

A fortiori, le \emph{$A$-torseur}
\[
\underline{\mathrm{Alb}}^{\chi(\delta)}_{X/S} \overset{A}{\times}
\underline{\mathrm{Pic}}^{\varphi(\delta)}_{B/S} \overset{A}{\times}
\bigl(P^{\delta} \overset{B}{\times} (A, \widetilde{\varphi(\delta)})\bigr)
\]
\emph{est trivial} (en fait avec une trivialisation canonique).

e) Pour expliciter ces considérations, il \struck{faudrait} en\-[core] \ill{}
avec pour tout $\delta \in \mathrm{NS}(X/S)$, « calculer » $\varphi(\delta)$
(et $\chi(\delta)$ \ldots). Pour ceci, on \struck{pas supposé} est ramené au cas
$S = \operatorname{Spec} k$, $k$ un corps, alg.\ clos si l'on veut.
Choisissons alors $p \in P^{\delta}(k)$, qui correspond : un faisceau inversible
$L$ sur $X$, dans la classe $\delta$.
\marginal{$X \times P^{\delta}$}
($(\mathcal{L}_g)_p = L$ \uncertain{sur} $X$.)
La donnée de $p$ définit un isom
\[ P^{\delta} \simeq B \]
et moyennant celui-ci, le faisceau $\mathcal{L}_g \mid X \times P^{\delta}$
\struck{\ill{}} se transforme en le faisceau
$\mathcal{L}^{0} \otimes \mathrm{pr}_1^{*}(L)$, où $\mathcal{L}^{0}$ est le
faisceau de Weil sur $X \times B$ [relatif au choix du point $g \in X(k)$] et
où $L$ est le faisceau inversible sur $X$ \struck{défini} de classe $\delta$
déjà envisagé. Il s'agit donc de calculer
\[
{\det}^{*}(R\,\mathrm{pr}_{2*})\bigl(\mathcal{L}^{0} \otimes \mathrm{pr}_1^{*}(L)\bigr)
\]

\page{6}

\struck{Si on fait} mod équivalence algébrique. \struck{Comme
$\mathrm{NS}(B/S)$ est \ill{} sans torsion, il suffit de faire} le calcul
\struck{(cela suffit si \ill{})} dans
$\mathrm{NS}(B/S) \otimes_{\mathbb{Z}} \mathbb{Q}$, \struck{et à fortiori
\ill{}}. On trouve donc les termes de degré 1 dans
$\mathrm{ch}_B\bigl(R\,\mathrm{pr}_{2*}(\mathcal{L} \otimes \mathrm{pr}_1^{*}(L))\bigr)$.
Or le $\mathrm{ch}_B$ se calcule par RR, du moins si $X$ est lisse, une
intersection complète, \ill{} on trouve
\[
\varphi(\delta) = \text{termes de degré 1 dans }
\mathrm{pr}_{2*}\bigl(\exp(D + \mathrm{pr}_1^{*}\Delta)\, \mathrm{pr}_1^{*}(\mathrm{Todd}_{X/k})\bigr)
\]
\[
= \mathrm{pr}_{2*}\bigl(\exp D \cdot \mathrm{pr}_1^{*}(T_{X/k} \exp \Delta)\bigr)
= [\exp D]\,(T_{X/k} \exp \Delta)
\]
\note{dans le second membre, un premier « $\mathrm{pr}$ » est surchargé à l'encre et « $\exp D$ » ajouté au crayon, entouré, au-dessus}
où
\[
\begin{cases}
D = \mathrm{cl}\,\mathcal{L} \in \mathrm{Gr}^{1}(X \times B) \\
\Delta = \mathrm{cl}\,\Delta \in \mathrm{Gr}^{1}(X) \\
\mathrm{Todd}_{X/k} \in \mathrm{Gr}^{*}(X)
\end{cases}
\]
\note{la page porte bien « $\Delta = \mathrm{cl}\,\Delta$ », là où l'on attendrait $\mathrm{cl}\,L$}

Nous allons faire le calcul en supposant $X$ lisse en dimension \ill{}, et
en travaillant en cohomologie $\ell$-adique ($\ell \cdot 1_k \neq 0$). On
choisit, pour simplifier, un isomorphisme $T_{\ell}(k^{*}) \simeq \mathbb{Z}_{\ell}$,
pour éviter le \uncertain{twist} de Tate. On a
\[
D \in H^{2}(X \times B) = H^{2}(X) + H^{1}(X) \otimes H^{1}(B) + H^{2}(B)
\]
et en fait
\[
D \in H^{1}(X) \otimes H^{1}(B)
\]
(et $D$ définit un isomorphisme entre $H^{1}(B)$ et $H^{1}(X)^{\vee}$),

\page{7}

On peut écrire
\[
D = \sum e_i \otimes e'_i
\]
où $(e_i)$ est une base de $H^{1}(X)$, et $e'_i$ \uncertain{est la base duale}
de $H^{1}(B)$. Posant
\[
\mathrm{Todd}_{X/k} \exp \Delta = \sum_{i=0}^{n} a^{i} \qquad a^{i} \in H^{2i}(X)
\]
on trouve
\[
\exp D\, \mathrm{pr}_1^{*}(\mathrm{Todd}_{X/k} \exp \Delta) = \sum_{i,j} \mathrm{pr}_1^{*}(a^{i})\, D^{j} / j!
\]
le terme en $i,j$ étant de degré cohomologique \uncertain{partiel} en $X$ égal
à $2i + j$. Il \uncertain{donne une} contribution nulle dans $\mathrm{pr}_{2*}$
\uncertain{sauf} si $2i + j = 2n$ i.e.\ $j = 2(n-i)$, \struck{donc il faut}
\add{\uncertain{auquel cas} on \ill{}} un $\mathrm{pr}_{2*}$ de degré $j$, qui
ne nous intéresse que si $j = 2$ i.e.\ $n - i = 1$. \struck{prendre} Donc on
trouve
\struck{$\sum \frac{1}{2}\, \mathrm{pr}_1^{*}(a_{n-1})\, D^{2}$ \ill{}}
\[
\varphi(\delta) = \tfrac{1}{2}\, \mathrm{pr}_{2*}\bigl(\mathrm{pr}_1^{*}(a_{n-1})\, D^{2}\bigr)
\]
\note{ici l'indice passe en bas : $a_{n-1}$ pour le $a^{n-1}$ de la ligne précédente}

\emph{Ex.\ 1} Prenons le cas $n = 1$. Alors $\delta$ est caractérisé par son
degré, soit \struck{\ill{}} $d$. On trouve
\struck{$\mathrm{Todd}_{X/k} \exp \Delta = (1 - \tfrac{K}{2})(1 + \Delta) = 1 + (\Delta - \tfrac{K}{2})$}
\marginal{\struck{degré} $d - g + 1$}\note{à droite de la ligne biffée, reliée à elle par un trait}
\[
\varphi(\delta) = \tfrac{1}{2}\, \mathrm{pr}_{2*}(D^{2})
\]
\struck{et il faut prendre $\mathrm{pr}_{2*}$ de $\frac{1}{2} D^{2}\, \mathrm{pr}_1^{*}(\Delta - \frac{K}{2})$}

Or \ill{} on a, si $(e_1, f_1, e_2, f_2, \ldots, e_g, f_g)$ est une base
symplectique de $H^{1}(X)$ [$e_i f_i = -f_i e_i = \eta$, base can.\ de
$H^{2}(X)$],
\marginal{on identifie $H^{1}(X)^{*}$ à \ill{} grâce à la \uncertain{dualisation}}
\[
D = \sum (e_i \otimes f_i + f_i \otimes e_i)
\]

\page{8}

donc
\[
D^{2} = \sum_{i,j} \underbrace{(e_i \otimes f_i + f_i \otimes e_i)(e_j \otimes f_j + f_j \otimes e_j)}
\]
\[
- [\, e_i f_j \otimes f_i e_j + f_i e_j \otimes e_i f_j \,]
\]
\[
= -\sum_i (\eta \otimes f_i e_i - \eta \otimes e_i f_i) = 2\eta \sum e_i f_i
\]
\note{la deuxième ligne, sous l'accolade, donne le terme général du produit ; la page ne l'indexe pas}
donc
\[
\tfrac{1}{2} D^{2} = \eta \otimes \sum e_i f_i
\]
et par suite
\[
\varphi(\delta) = \sum e_i \wedge f_i \in H^{2}(B)
\]
\emph{C'est la classe de polarisation canonique de la jacobienne $B$.}
D'ailleurs
\[
\chi(\delta) = 1 - g + d
\]
\marginal{vérifier que le signe \uncertain{est} correct}
\struck{\ill{}} Utilisons la polarisation canonique \add{par le diviseur
$\Theta$} de $B$, définissant un isom can
\[
c : B \longrightarrow A
\]
et transformons les $A$-torseurs de la fin de d) par $c^{-1}$. On sait
\emph{que $\underline{\mathrm{Alb}}^{1}_{X/S}$ se transforme en}
\[
\underline{\mathrm{Pic}}^{-1}_{X/S} = P^{-1}
\quad \text{(\emph{attention au signe !})}
\]
donc $\underline{\mathrm{Alb}}^{\chi(\delta)}_{X/S}$ devient
$P^{-\chi(\delta)} = P^{g-1-d}$, d'autre part
$\underline{\mathrm{Pic}}^{\varphi(\delta)}_{B/S}$
\add{$\varphi(\delta) = $ \uncertain{pol.\ can.\ de $A$}} \struck{devient}
\struck{$\underline{\mathrm{Pic}}^{1}$} $P^{g-1}$ (isom.\ canonique défini par
le div.\ $\Theta$), enfin
\[
P^{\delta} \overset{B}{\times} (A, \underset{\substack{\parallel \\ c}}{\widetilde{\varphi(\delta)}})
\;\text{devient}\; P^{\delta} = P^{d} .
\]
Donc le torseur envisagé (produit des trois précédents) est $P^{2g-2}$. On
trouve \ill{} une section canonique de \ill{} : ce n'est \uncertain{pas très
excitant}, car on s'attend \uncertain{à trouver} la section canonique $K$
(à prouver !).

\page{9}

\emph{N.B.} Pour bien faire, il faudrait prouver que sur le \uncertain{schéma}
modulaire $S$ des schémas en courbes \add{\ill{}} de genre $g$
\struck{$\geq 1$} ($g \geq 1$ ; on \struck{\ill{}} \uncertain{choisit} une
section marquée pour $g = 1$), \uncertain{tel que} si $\mathcal{X}$ est la
courbe de genre $g$ \struck{\ill{}} universelle, $P = \underline{\mathrm{Pic}}_{X/S}$,
alors le groupe des sections de $P$ sur $S$ est engendré par $K$
(\struck{\ill{}} \add{en particulier, que} le groupe des sections de $P^{0}$
sur $S$ est égal à zéro). Il suffirait de déterminer le groupe en question
\uncertain{en car.\ nulle}.
\note{l'insertion interlinéaire au-dessus de « courbes », longue d'une ligne, est illisible}

f) \emph{Ex.\ 2}\quad \struck{\ill{}} \emph{Cas} $X$ un schéma abélien. Donc
$X = A$. On sait que \add{sur les \uncertain{cycles} \ill{}, $\exp D$
\uncertain{établit}} \struck{\ill{}}
\[
\Phi : \mathrm{chow}(A)_{\mathbb{Q}} \xrightarrow{\ \sim\ } \mathrm{chow}(B)_{\mathbb{Q}}
\]
qui \struck{\ill{}} transforme élément de degré $i$ en élément de degré
complémentaire $n - i$, \add{et} qui transforme produit ordinaire en produit de
convolution, et inversement. D'ailleurs $T_{X/k} = 1$, donc on trouve
\add{composante de degré 1 de}
\[
\varphi(\delta) = \sum \frac{1}{i!}\, \underbrace{\Phi(\Delta)}_{\text{classe d'une courbe}}{}^{*i}
= \frac{1}{(n-1)!}\, \Phi(\Delta)^{*(n-1)} .
\]
\struck{donc} On trouve ainsi un hom (\add{\uncertain{polynomial}} de degré
$n - 1$ \ill{})
\[
\varphi : \underline{\mathrm{NS}}_{A/S} \longrightarrow \underline{\mathrm{NS}}_{B/S}
\]
On vérifie que, si $\delta$ est une section de $\underline{\mathrm{NS}}_{A/S}$,
définissant
\[ \tilde{\delta} : A \longrightarrow B \]
et $\varphi(\delta)$ section de $\underline{\mathrm{NS}}_{B/S}$, d'où
\[ \widetilde{\varphi(\delta)} : B \longrightarrow A \]
\marginal{Faire calcul direct, en faisant \struck{\ill{}} attention aux signes !}\note{note entourée, en bas à gauche}
\[
\boxed{\widetilde{\varphi(\delta)}\, \tilde{\delta} = -\chi(\delta)\, \mathrm{id}_A ,
\qquad \tilde{\delta}\, \widetilde{\varphi(\delta)} = -\chi(\delta)\, \mathrm{id}_B}
\]
\note{le signe moins du second membre est moins net que celui du premier}

\page{10}

Par suite, si $\delta$ est non dégénérée i.e.\ $\chi(\delta) \neq 0$, alors
$\varphi(\delta)$ est non dégénérée, \struck{et si $\delta$ est une
polarisation,} \add{\uncertain{déduit}}
$-\varphi(\delta) = \delta'$ est une polarisation. On \add{\ill{}}
d'ailleurs de la relation plus haut, puisque
$\deg \tilde{\delta} = \chi(\delta)^{2}$,
$\deg \tilde{\delta}' = \chi(\delta')^{2}$,
$\deg \tilde{\delta}\tilde{\delta}' = \chi(\delta)^{2}\chi(\delta')^{2}
= \chi(\delta)^{2n}$ :
\[
\chi(\delta') = \chi(\delta)^{n-1} \qquad \delta' = -\varphi(\delta)
\]
\note{dans l'égalité encadrée des degrés, la dernière valeur $\chi(\delta)^{2n}$ est écrite sous le produit, reliée par deux traits verticaux ; l'exposant $n-1$ est surchargé}
\marginal{$\delta' = -\varphi(\delta)$ ; \uncertain{je pense}, \ill{} \ill{} au carré ; \ill{} \ill{} dimension \ill{} \uncertain{prouver} \ill{} si $\delta$ est une polarisation}\note{note oblique dans la marge gauche, en partie illisible}
si $\varphi(\delta)$ est non dégénérée, \struck{\ill{} $\delta$ \ill{}}
$\chi(\delta)$ l'est.

Notons aussi que si $\delta' = \varphi(\delta)$, $\delta'' = -\varphi(\delta')$, on a
\[
\tilde{\delta}' = \tilde{\delta}^{-1} \chi(\delta) ,
\qquad
\tilde{\delta}''^{-1} = \tilde{\delta}'^{-1} \chi(\delta')
= \tilde{\delta}\, \chi(\delta)^{-1} \chi(\delta') = \tilde{\delta}\, \chi(\delta)^{n-2}
\]
\note{devant $\tilde{\delta}''^{-1}$, un premier « $\varphi(\varphi(\delta))$ » est biffé}
donc
\[
\varphi(\varphi(\delta)) = \delta\, \chi(\delta)^{n-2}
\]
\note{les signes sont ceux de la page ; aucun n'y est ajusté}
Si \struck{\ill{}} $n = 1$, i.e.\ $X$ une courbe elliptique, on
\struck{$\varphi(\delta) = $ polarisation canonique,} trouve $\delta\,\chi(\delta)$
i.e.\ $\varphi(\varphi(\delta))$ est la polarisation canonique ;
\struck{[N.B.\ $\varphi($} mieux, pour tout $\delta$, $\varphi(\delta)$ est la
polarisation canonique. Si $n \geq 2$, la formule
$\varphi(\varphi(\delta)) = \delta\,\chi(\delta)^{n-2}$ inspire \uncertain{confiance}
directement\ldots

\emph{N.B.} Il faudrait prouver que sur le \struck{schéma} \struck{\ill{}}
\add{site} modulaire des schémas abéliens polarisés $(A, \delta)$
\add{\uncertain{en car.\ 0}}, avec degré de polarisation donné $\chi(\delta)$,
la polarisation $\varphi(\delta) = \delta'$ \struck{\ill{}} sur $B$ n'est plus
divisible \struck{i.e.\ que \ill{}} \ill{} ce qui montrerait que $\varphi(\delta)$
est ce qu'on peut faire de mieux dans ce genre, comme \add{\uncertain{bien
fait}} « inversion » d'une polarisation. \struck{On} On s'attend à ce que
$\varphi(\delta)$ engendre le groupe \struck{\ill{} de \ill{}} des sections de
$\underline{\mathrm{NS}}_{B/S}$ sur une quelconque \struck{\ill{}} site
modulaire\ldots\ \struck{\ill{}} \ill{} comprenant comme \ill{} \struck{site
modulaire}

\page{11}

L'isomorphisme de la fin de d) nous fournit une section canonique de
\[
\underline{\mathrm{Pic}}^{-\delta'}_{B/S} \overset{A}{\times}
\bigl(\underline{\mathrm{Pic}}^{\delta}_{A/S} \overset{B}{\times} (A, \tilde{\delta}')\bigr)
\]
où $\delta' = -\varphi(\delta)$ comme dessus. En d'autres termes, on trouve un
isomorphisme canonique
\[
(*) \qquad
\boxed{\underline{\mathrm{Pic}}^{\delta'}_{B/S} \simeq
\underline{\mathrm{Pic}}^{\delta}_{A/S} \overset{B}{\times} (A, \tilde{\delta}')}
\]
\note{dans le cadre, l'exposant $\delta$ de $\underline{\mathrm{Pic}}_{A/S}$ est écrit par-dessus un premier exposant, peut-être $-\delta$ ; à droite du cadre, au crayon, « ${}_{\mathbb{Z}}A$ », lecture douteuse}

\marginal{N.B.\ Comme $\delta$ \uncertain{lisse}, $\underline{\mathrm{Pic}}^{-\delta}_{A/S} = \underline{\mathrm{Pic}}^{\delta}_{A/S}$ (canoniquement) on peut \uncertain{supprimer} \ill{} \uncertain{signes} dans \ill{}\ldots\ Prouvons que \uncertain{pour} $\delta$ \uncertain{ample} \ill{} \uncertain{torseurs} ? \ill{} un isomorphisme \ill{} \uncertain{compatible} avec \ill{} \uncertain{canoniques}, \ill{} \uncertain{section} $(*)$ \ill{} \ill{} \ill{} \ill{} \ill{} \ill{} \ill{} $A$\ldots\ à vérifier}\note{longue note oblique dans la marge gauche, cernée d'un trait, en grande partie illisible}

\emph{Exemple} Supposons que $\delta$ soit de degré 1, \struck{\ill{}}
$\tilde{\delta} : A \xrightarrow{\sim} B$, et $\tilde{\delta}' : B \xrightarrow{\sim} A$
est l'inverse. Alors $\delta'$ \struck{\ill{}} \add{est la} polarisation
déduite de $\delta$ par transport de structure : l'inverse de
$\tilde{\delta}$, $\underline{\mathrm{Pic}}^{\delta'}_{B/S} \simeq
\underline{\mathrm{Pic}}^{\delta}_{A/S}$ [\ill{} on trouve un isom., via
$A \xrightarrow{\tilde{\delta}} B$] et l'isom.\ \add{de $B$-torseurs}
(compatible \struck{\ill{}} avec \ill{}) \uncertain{canonique} peut s'interpréter
comme un isom.
\[
\underline{\mathrm{Pic}}^{\delta}_{A/S} \simeq \underline{\mathrm{Pic}}^{-\delta}_{A/S}
\]
\struck{\ill{} bien connu, \ill{} l'isomorphisme canonique \ill{} défini par
la \uncertain{symétrie} \ill{} $(-1)$ \ill{} $(-1)$ \ill{} et en diviseurs \ill{}}
\marginal{\struck{$A \times B$}}
\note{un bloc de quatre lignes est encadré et biffé de grandes boucles, et une ligne suivante est raturée d'une hachure serrée ; on n'en lit que des bribes}

\emph{Remarque} La section du torseur de la fin de d) est stable par
automorphismes de $(X/S, \delta)$ dans le cas actuel, \ill{} par
\struck{\ill{}} automorphismes par translations dans $X = A$. \ill{} Il ne peut
pas \uncertain{venir} \ill{} de section par ce qui suit. On trouve
\[
\underline{\mathrm{Alb}}^{\chi(\delta)}_{X/S} \times
\underline{\mathrm{Pic}}^{\varphi(\delta)}_{B/S} \overset{A}{\times}
\bigl(P^{\delta} \overset{B}{\times} (A, \widetilde{\varphi(\delta)})\bigr)
= \underline{\mathrm{Pic}}^{\varphi(\delta)}_{B/S} \overset{A}{\times}
\bigl(P^{\delta} \overset{B}{\times} (A, \widetilde{\varphi(\delta)})\bigr)
\]
par translation par \ill{}
$\chi(\delta)\, a + \widetilde{\varphi(\delta)}\, \tilde{\delta}(a)$ \struck{\ill{}} ;
donc on a $[\chi(\delta)\, \mathrm{id}_A + \widetilde{\varphi(\delta)}\, \tilde{\delta}]\, a = 0$
pour tout $a$, d'où \struck{\ill{}}
$\widetilde{\varphi(\delta)}\, \tilde{\delta} = -\chi(\delta)\, \mathrm{id}_A$.
\marginal{Attention aux signes ??}\note{en bas à gauche, cerné au crayon}

\page{12}

\emph{Remarque} Dans l'isom can $(*)$, remplaçons $\delta$ par $\lambda\delta$
($\lambda \in \mathbb{Z}$). Les deux membres sont remplacés respectivement par
\[
\underline{\mathrm{Pic}}^{\lambda^{n-1}\delta'}_{B/S} \simeq
\bigl(\underline{\mathrm{Pic}}^{\delta'}_{B/S}\bigr)^{(\lambda^{n-1})}
\]
et par
\[
\underline{\mathrm{Pic}}^{-\lambda\delta}_{A/S} \overset{B}{\times} (A, \lambda^{n-1}\tilde{\delta}')
\simeq \bigl(\underline{\mathrm{Pic}}^{-\delta}_{A/S} \overset{B}{\times} (A, \tilde{\delta}')\bigr)^{(\lambda^{n})} .
\]
\note{les exposants en $\lambda$ sont tracés très vite ; on lit $\lambda^{n-1}$ et $\lambda^{n}$ d'après la ligne suivante, sans certitude}
Donc on trouve, \uncertain{écrivant} l'égalité des deux membres pour
$\lambda\delta$, que
\[
\bigl(\underline{\mathrm{Pic}}^{\delta'}_{B/S}\bigr)^{\otimes(\lambda^{n} - \lambda^{n-1})}
\simeq \text{torseur trivial (canonique)}
\]
Faisant $\lambda = -1$, on trouve
\[
\bigl(\underline{\mathrm{Pic}}^{\delta'}_{B/S}\bigr)^{\otimes 2}
\simeq \text{torseur trivial} \quad \text{(canonique)}
\]
[\emph{N.B.} 2 est le \uncertain{p.g.c.d.} \uncertain{des valeurs} de la forme
$\lambda^{n} - \lambda^{n-1}$].
\marginal{Vérifier que la trivialisation \add{et la trivialisation} canonique
\uncertain{définie} pour \struck{\ill{}} \uncertain{toute} section $\delta'$ de
$\underline{\mathrm{NS}}_{B/S}$}\note{note oblique dans la marge gauche}

\struck{Lorsque} Comme $(\delta')' = \chi(\delta)^{n-2}\delta$, on en conclut,
en échangeant les rôles de $A$ et $B$
\[
\bigl(\underline{\mathrm{Pic}}^{\delta}_{A/S}\bigr)^{\otimes(2\chi(\delta)^{n-2})}
\simeq \text{torseur trivial (canon !)}
\]
\marginal{si $n \geq 2$ (trivial si $n \leq 1$)}
Lorsque $\chi(\delta)$ est pair, on a même
\[
\bigl(\underline{\mathrm{Pic}}^{\delta}_{A/S}\bigr)^{\otimes(\chi(\delta)^{n-2})}
\simeq \text{torseur trivial (canonique)}
\]
pour $\chi(\delta)$ \emph{pair}, \uncertain{si $n \geq 2$}, \ill{} \ill{} pas
\ill{} les \uncertain{espérances} qui suivent\ldots
\marginal{si $n \geq 3$, vérifier que cette section est \uncertain{encore} la
section canonique de
$\bigl(\underline{\mathrm{Pic}}^{\delta\chi(\delta)^{n-2}/2}_{A/S}\bigr)^{\otimes 2}$,
\ill{} $\chi(\delta)$ pair, \ill{} et que \ill{} $\delta$ divisible par 2,
\ill{} il faudrait \uncertain{prouver} \ill{} que la trivialisation \ill{} la
trivialisation canonique}\note{longue note oblique dans la marge gauche, en bas}

\struck{Lorsque $\chi(\delta) = 1$, et que $\delta$ est une polarisation, il
faudrait \ill{} les sections du $(\underline{\mathrm{Pic}}^{\delta}_{A/S})^{\otimes 2}$
\uncertain{obtenues} \ill{} \ill{} l'isom
$\underline{\mathrm{Pic}}^{\delta} \simeq \underline{\mathrm{Pic}}^{-\delta}$
transformant \ill{} (1) \ill{} $(-1)$ \ill{}}
\note{le dernier paragraphe est encadré et biffé de trois grandes diagonales ; les nombres $(1)$ et $(-1)$ y sont entourés}

\page{13}

\struck{g) Généralisation des constructions précédentes\ldots}

\struck{\ill{}} (Soit $\delta$ \struck{polarisation \ill{}} \add{section} de
$\underline{\mathrm{NS}}_{A/S}$, $\delta' = \ldots$ $\delta'$
\struck{polarisation} \add{\uncertain{de}} $B$, $\delta''$
\struck{polarisation de} $A$.)
\note{la première ligne est biffée ; « Remarque » (lecture douteuse) au-dessous est surchargé, et la parenthèse qui suit reste ouverte sur la page}
On a successivement
\[
\underline{\mathrm{Pic}}^{\delta'}_{B/S} \simeq
\underline{\mathrm{Pic}}^{\delta}_{A/S} \overset{B}{\times} (A, \tilde{\delta}')
\]
\[
\underline{\mathrm{Pic}}^{\delta''}_{A/S} \simeq
\underline{\mathrm{Pic}}^{\delta'}_{B/S} \overset{A}{\times} (B, \tilde{\delta}'')
\simeq \underline{\mathrm{Pic}}^{\delta}_{A/S} \overset{B}{\times} (B, \tilde{\delta}''\tilde{\delta}')
\]
or on a
\[
\begin{cases}
\delta'' = \delta\, \chi(\delta)^{n-2} \\
\tilde{\delta}''\tilde{\delta}' = \chi(\delta')\, \mathrm{id}_B = \chi(\delta)^{n-1}\, \mathrm{id}_B
\end{cases}
\]
donc les \ill{} colonnes de \uncertain{ces} \ill{}
\[
\underline{\mathrm{Pic}}^{\delta\chi(\delta)^{n-2}}_{B/S} \simeq
\underline{\mathrm{Pic}}^{\delta\chi(\delta)^{n-1}}_{B/S}
\]
d'où
\[
\bigl(\underline{\mathrm{Pic}}^{\delta}_{B/S}\bigr)^{\otimes[\chi(\delta)^{n-1} - \chi(\delta)^{n-2}]}
\simeq \text{torseur trivial (isom.\ canonique)}
\]
\note{ici $\underline{\mathrm{Pic}}^{\delta}$ porte l'indice $B/S$, là où les deux lignes précédentes mettaient $A/S$ ; on transcrit tel quel}
Remplaçant $\delta$ par $\lambda\delta$ ($\lambda \in \mathbb{Z}$), on trouve
plus généralement
\[
\bigl(\underline{\mathrm{Pic}}^{\delta}_{B/S}\bigr)^{\otimes \chi(\delta)^{n-2}(\lambda^{n}\chi(\delta) - \lambda^{n-1})}
\simeq \text{torseur trivial}
\]
\note{l'exposant est surchargé : un premier facteur en tête est raturé, et « $\lambda^{n}\chi(\delta) - \lambda^{n-1}$ » est repris par-dessus une première écriture}
\struck{$\lambda^{n-1}(\lambda\chi(\ldots)$ \ill{} $\chi(\delta)^{n-2}[\lambda^{n}\chi(\delta) - \lambda^{n-1}]$}
\struck{Soit $\alpha$ le p.g.c.d.} \uncertain{p.p.c.m.} des
$\lambda^{n}\chi(\delta) - \lambda^{n-1}$ ($\lambda \in \mathbb{Z}$) :
\uncertain{donc} \ill{} $p$ divise $\alpha$ \uncertain{ssi} \ill{}
$\lambda^{n}\chi(\delta) - \lambda^{n-1} \equiv 0 \ (p)$ pour tout $\lambda \in \mathbb{Z}$
\note{le passage « Soit $\alpha$ \ldots\ pour tout $\lambda \in \mathbb{Z}$ » est cerné d'un trait et barré en partie ; la lecture de « p.p.c.m. » et du rôle de $p$ est douteuse}
Faisant $\lambda = \pm 1$ et \uncertain{ajoutant}, on \uncertain{voit} \ill{}
\[
\bigl(\underline{\mathrm{Pic}}^{\delta}_{B/S}\bigr)^{\otimes 2\chi(\delta)^{n-2}}
\simeq \text{torseur trivial}
\]
et \uncertain{de même}
\[
\bigl(\underline{\mathrm{Pic}}^{\delta}_{B/S}\bigr)^{\otimes \chi(\delta)^{n-2}}
\simeq \text{torseur trivial si $\chi(\delta)$ pair}
\]
[\uncertain{car alors} $\chi(\delta) - 1$ est impair, et on utilise que
\[
\bigl(\underline{\mathrm{Pic}}^{\delta}_{B/S}\bigr)^{(\chi(\delta)-1)\chi(\delta)^{n-2}}
\simeq \text{torseur trivial}.
\]
\emph{Exemple} \uncertain{si} $n = 2$, \ill{} $\underline{\mathrm{Pic}}^{\delta}$
trivial \uncertain{et} $(\underline{\mathrm{Pic}}^{\delta}_{B/S})^{\otimes 2}$
trivial, si $\chi(\delta)$ pair\ldots

\page{14}

e) Généralisation des constructions précédentes.
\note{la page reprend la lettre e), déjà prise p.~5 ; la rubrique « g) Généralisation\ldots » biffée en tête de la page 13 annonçait ce paragraphe}

On revient aux conditions de a), mais on suppose donné un complexe
\struck{parfait} $\mathcal{F}^{\bullet}$ sur $X$, \struck{\ill{} complexe}
parfait (ou seulement parfait rel$/S$) et on considère
\[
{\det}^{*} R\,\mathrm{pr}_{2*}\bigl(\mathrm{pr}_1^{*}(\mathcal{F}^{\bullet}) \otimes \mathcal{L}_g\bigr) = M_g
\]
On aura, pour une autre choix $g'$, \struck{et} un isom.\ can.
\[
M_{g'} \simeq M_g\, g'^{*}_{P}(\mathcal{L}_g)^{-\chi_{\mathcal{F}}}
= M_g\, (N_g N_{g'}^{-1})^{\chi_{\mathcal{F}}}
\]
\ill{} maintenant $\chi$ est la fonction \uncertain{rang} \ill{} pour
$R\,\mathrm{pr}_{2*}(\mathrm{pr}_1^{*}(\mathcal{F}) \otimes \mathcal{L}_g)$.
On trouve alors, pour $\mathcal{F}$ fixe, un homomorphisme
\[
\varphi_{\mathcal{F}^{\bullet}} : \underline{\mathrm{NS}}_{X/S} \longrightarrow \underline{\mathrm{NS}}_{B/S}
\]
dépendant additivement de $\mathcal{F}^{\bullet}$ (pour la structure de
catégorie \emph{triangulée} \uncertain{de} $D_{\mathrm{parf}}(X)$ resp.\
$D_{\mathrm{parf}}(X/S)$). On trouve alors, pour toute section $\delta$ de
$\underline{\mathrm{NS}}_{X/S}$, un hom
\[
X \xrightarrow{\ \ell^{\mathcal{F}^{\bullet}}_{\delta}\ }
\underline{\mathrm{Pic}}^{\varphi_{\mathcal{F}}(\delta)}_{B/S} \overset{A}{\times}
\bigl(P^{\delta} \overset{B}{\times} (A, \widetilde{\varphi_{\mathcal{F}}(\delta)})\bigr)
\]
d'où un homom.\ de $A$-torseurs
\[
\underline{\mathrm{Alb}}^{-\chi_{\mathcal{F}}(\delta)}_{X/S} \longrightarrow
\underline{\mathrm{Pic}}^{\varphi_{\mathcal{F}}(\delta)}_{B/S} \times
\bigl(\underline{\mathrm{Pic}}^{\delta}_{X/S} \overset{B}{\times} (A, \widetilde{\varphi_{\mathcal{F}}(\delta)})\bigr)
\]
et enfin une section
\[
\boxed{L_{\mathcal{F}}(\delta) \in \Gamma\Bigl(S, \underline{\mathrm{Alb}}^{\chi_{\mathcal{F}}(\delta)}_{X/S}
\overset{A}{\times} \underline{\mathrm{Pic}}^{\varphi_{\mathcal{F}}(\delta)}_{B/S}
\times \bigl(\underline{\mathrm{Pic}}^{\delta}_{X/S} \overset{B}{\times} (A, \widetilde{\varphi_{\mathcal{F}}(\delta)})\bigr)\Bigr)}
\]
\note{dans les deux dernières formules, le troisième facteur est écrit $\underline{\mathrm{Pic}}^{\delta}_{X/S}$ là où la p.~4 avait $P^{\delta}$ ; le signe de l'exposant de $\underline{\mathrm{Alb}}$ diffère d'une ligne à l'autre, comme p.~4}

Il faut enfin calculer $\chi_{\mathcal{F}}(\delta)$ et
$\varphi_{\mathcal{F}}(\delta)$, et pour \ill{} \ill{} on est ramené \ill{}
\uncertain{au} cas, comme dans e), où $S = \operatorname{Spec} k$, $k$ corps
alg.\ clos, et à calculer
\[
\chi(L_{\delta} \otimes \mathcal{F}^{\bullet}) \quad\text{et}\quad
{\det}^{*} R\,\mathrm{pr}_{2*}\bigl(\mathcal{L} \otimes \mathrm{pr}_1^{*}(L_{\delta} \otimes \mathcal{F})\bigr)
\]
On note que dans le calcul, $\delta$ disparaît, \add{$\mathcal{F}$ \uncertain{près}}
\uncertain{à condition} de remplacer \struck{$\mathcal{F}$} $\mathcal{F}$ par
$L_{\delta} \otimes \mathcal{F}$\ldots

\page{15}

On peut le faire par RR si $X$ est loc.\ une intersection complète (rel$/S$),
on trouve
\struck{$\chi(L_\delta)$}
\[
\chi_{\mathcal{F}}(\delta) = \chi(L_{\delta} \otimes \mathcal{F}^{\bullet})
= \pi\bigl(\mathrm{ch}(L_{\delta} \otimes \mathcal{F})\, \mathrm{Todd}(X)\bigr)
\]
\[
\varphi_{\mathcal{F}}(\delta) = \text{terme de degré 1 dans }
(\exp D) . [\mathrm{ch}_X(L_{\delta} \otimes \mathcal{F})\, \mathrm{Todd}(X)]
\]
\note{la lettre devant la parenthèse, dans la première formule, est lue $\pi$ ; à droite, une ligne ajoutée précise : « $=$ terme constant dans $\mathrm{pr}_{2*}\bigl(\mathrm{pr}_1^{*}(\mathrm{ch}_X(L_{\delta} \otimes \mathcal{F}))\,\mathrm{Todd}(X)\, \exp_{X \times B}(D)\bigr)$ », avec sous l'accolade « $(\exp D).[\mathrm{ch}_X(L_{\delta} \otimes \mathcal{F})\,\mathrm{Todd}(X)]$ » ; « $\mathrm{Todd}(X)$ » y est entouré et inséré}

\emph{Ex.\ 1} $X$ \uncertain{courbe}$/S$ de dim.\ relative 1. On aura
\[
\mathrm{ch}_X(L_{\delta} \otimes \mathcal{F}) =
(1 + \underset{\substack{\mid \\ c(L_{\delta})}}{\delta})
(\underset{\substack{\parallel \\ \mathrm{rang}\,\mathcal{F}^{\bullet}}}{\rho} + \underset{\substack{\parallel \\ \det'(\mathcal{F})}}{\gamma})
= 1 + (\delta + \gamma)
\]
\note{sous le premier membre, « $\delta$ \uncertain{de degré} $d$ » ; sous $\delta + \gamma$, « degré $d$ » et « degré $c$ » ; le $1$ du second membre est écrit ainsi bien que le premier facteur donne $\rho$. Les lettres $\rho$ et $\gamma$ sont lues sans certitude}
\[
\mathrm{Todd}(X) = 1 + \underset{\text{degré } 1-g}{\tfrac{1}{2}K}
\]
\[
\mathrm{ch}_X(L_{\delta} \otimes \mathcal{F}) \otimes \mathrm{Todd}(X) = 1 + \tau ,
\qquad \text{avec } \deg \tau = 1 - g + d + c
\]
\struck{Il faut prendre les \ill{}}
\[
\boxed{\chi_{\mathcal{F}}(\delta) = 1 - g + d + c}
\]
D'autre part, $\varphi(\delta)$ ne \uncertain{peut} changer \add{\uncertain{ici}}
\struck{\ill{}} par rapport au calcul précédent dans e)
\[
\boxed{\varphi_{\mathcal{F}}(\delta) = \varphi(\delta) = \tfrac{1}{2}\, \mathrm{pr}_{2*}(D^{2})
= \text{polarisation canonique de } B}
\]
Donc on trouve une section canonique de
\[
\underline{\mathrm{Alb}}^{\chi_{\mathcal{F}}(\delta) = \chi(\delta) + c}_{X/S}
\overset{A}{\times} \underline{\mathrm{Pic}}^{\varphi(\delta)}_{B/S}
\overset{A}{\times} \bigl(\underline{\mathrm{Pic}}^{\delta}_{X/S} \overset{B}{\times} (B, \widetilde{\varphi(\delta)})\bigr)
\]
\note{dans le dernier facteur, l'exposant de $\underline{\mathrm{Pic}}_{X/S}$ est surchargé et la lettre sous le produit est peu nette ; la page écrit $(B, \ldots)$, là où la p.~4 avait $(A, \ldots)$}
donc, en utilisant le calcul fait dans e), on trouve une section canonique de
$\underline{\mathrm{Alb}}^{c}_{X/S} \simeq \underline{\mathrm{Pic}}^{c}_{X/S}$.
Il n'y a pas de miracle ici, car il suffit de prendre la section définie par
\add{le \uncertain{produit} \uncertain{inversible}} $\det^{*}\mathcal{F}$ sur $X/S$.

\page{16}

\emph{Ex.\ 2} $X$ un schéma abélien $A$. On a ici $\mathrm{Todd}(X) = 1$.
On trouve
\[
\chi_{\mathcal{F}}(\delta) = \pi\bigl(\mathrm{ch}(L_{\delta}\mathcal{F})\bigr) = \deg a_n
\]
si on pose
\[
\mathrm{ch}(L_{\delta}\mathcal{F}) = \sum_{1}^{n} a_i \qquad a_i \in \mathrm{Chow}(A)
\]
On trouve ensuite, comme \uncertain{avant} \add{\uncertain{terme de degré 1 de}}
\[
\varphi_{\mathcal{F}}(\delta) = (\exp D) . (a_{n-1}) = \mathrm{pr}_{2*}\bigl(\tfrac{1}{2} D^{2} a_{n-1}\bigr) .
\]

f) Nous allons donner une autre interprétation de l'homomorphisme $\ell_{\delta}$
de d), en définissant \add{(en intégrant exclusivement sur $X \times_S B$)} un
homomorphisme
\[
\boxed{\ell' : X \times \underline{\mathrm{Pic}}_{X/S} \longrightarrow \underline{\mathrm{Pic}}_{B/S}}
\quad \text{(additif en $\underline{\mathrm{Pic}}_{X/S}$)} .
\]
\note{dans le cadre, l'indice de $\underline{\mathrm{Pic}}$ est écrit $X/B$ ou $X/S$ ; on lit $X/S$, comme dans la parenthèse}
Pour ceci, soit $g$ une section de $X$, et $\lambda$ une section de
$\underline{\mathrm{Pic}}_{X/S}$, qui (grâce à $g$) est définie par un module
inversible $L$ \add{$= L(\lambda, g)$} sur \struck{\ill{}} $X$, trivialisé le
long de $g$. On a par \uncertain{ailleurs} sur $X \times_S B$
($W_g = W(g)$ faisceau de Weil sur $X \times_S B$, \uncertain{correspondant à
$g$}), un \struck{\ill{}} faisceau inv.\ sur $B$,
\[
{\det}^{*}\bigl(\mathrm{pr}_{2*}(\mathrm{pr}_1^{*}(L) \otimes W)\bigr)
\]
\uncertain{définissant} une section $\ell(g, \lambda)$ dont la classe
\ill{} $\underline{\mathrm{Pic}}_{B/S}$. La formation est compatible avec
extension de la base, d'où l'hom $\ell$. Désignons, pour une section $g$ de
$X$, par $\ell_g$ l'hom $\underline{\mathrm{Pic}}_{X/S} \to \underline{\mathrm{Pic}}_{B/S}$
qu'elle définit.

\page{17}

On va préciser comment $\ell_g$ varie avec $g$. Soit $g'$ une autre section.
Alors on a
\[
L(\lambda, g') = L(\lambda, g)\, g'^{*}(L(\lambda, g))^{-1} =
\]
\struck{donc}
\[
W(g') = W(g)\, g'^{*}_{B}(W(g))^{-1}
\]
(où $g'_B$ est la section de $X \times_S B / B$ déduite de la section $g'$ de
$X/S$). On a donc
\[
\mathrm{pr}_1^{*}(L(\lambda, g'))\, W(g') = \mathrm{pr}_1^{*}(L(\lambda, g))\, W(g)\, \mathcal{D}
\qquad
\mathcal{D} = g'^{*}_{B}(L(\lambda, g))^{-1} \otimes_{\mathcal{O}_X} g'^{*}_{B}(W(g))^{-1}
\]
d'où comme d'habitude
\[
\ell_{g'}(\lambda) = \ell_g(\lambda) + \bigl\lbrace \text{classe dans } \underline{\mathrm{Pic}}_{B/S}
\text{ de } g'^{*}_{B}(W(g))^{-\chi(\lambda)} = (N_{g'} N_g^{-1})^{-\chi(\lambda)}
\]
\[
(*) \qquad = \ell_g(\lambda) + (\psi_X(g') - \psi_X(g))(-\chi(\lambda))
\]
où $\psi_X : X \to \underline{\mathrm{Alb}}^{1}_{X/S}$ est l'hom canonique, et
$\underline{\mathrm{Alb}}^{1}_{X/S}$ est considéré comme torseur sous
$A = \underline{\mathrm{Pic}}^{0}_{B/S}$ de la façon habituelle. Donc,
considérant l'hom
\[
\varphi : \underline{\mathrm{NS}}_{X/S} \longrightarrow \underline{\mathrm{NS}}_{B/S}
\]
\uncertain{induit} par $\ell$ de façon évidente, on déduit pour toute section
$\delta$ de $\underline{\mathrm{NS}}_{X/S}$ un hom
\[
\boxed{\ell_{\delta} : X \times \underline{\mathrm{Pic}}^{\delta}_{X/S} \longrightarrow \underline{\mathrm{Pic}}^{\varphi(\delta)}_{B/S}}
\]
\struck{qui, \ill{} définit \ill{} \ill{} \ill{} factorisation $\underline{\mathrm{Alb}}^{\chi(\delta)}_{X/S}$ d'ailleurs par $\underline{\mathrm{Alb}}^{-\chi(\delta)}_{X/S} \overset{A}{\times}$ \ill{} satisfaisant \ill{}}
\note{les deux dernières lignes sont biffées d'un trait chacune}

\page{18}

De plus, pour une section $g$ de $X/S$, on vérifie que l'on a, pour
$\beta \in \underline{\mathrm{Pic}}^{0}_{X/S}(S) = B(S)$,
\[
(**) \qquad \ell_g(\lambda + \beta) = \ell_g(\lambda) + \widetilde{\varphi(\delta)}(\beta)
\qquad \bigl(\lambda \in \underline{\mathrm{Pic}}^{\delta}_{X/S}(S)\bigr)
\]
\note{le signe entre $\ell_g(\lambda)$ et $\widetilde{\varphi(\delta)}(\beta)$ est noyé sous une tache d'encre ; on lit $+$ d'après le reste de la page}

\struck{En effet, on sait que (comme toute flèche \uncertain{entre} torseurs
sous des schémas abéliens \uncertain{au-dessus de} $S$)
$\ell^{\delta}_g : \underline{\mathrm{Pic}}^{\delta}_{X/S} \to \underline{\mathrm{Pic}}^{\varphi(\delta)}_{B/S}$
est un morphisme [de] torseurs sous $B$ [dans] torseurs sous $A$ \ill{} un
hom $\rho : B \to A$, et il faut prouver que $\rho = -\widetilde{\varphi(\delta)}$.
On \uncertain{peut} pour ceci supposer $S$ le spectre d'un corps alg.\ clos $k$,
et il suffit de prouver que l'on a \uncertain{égalité} pour \ill{}
\uncertain{éléments} \ill{} $B(k)$, i.e.\ \ill{} [\uncertain{les}] \ill{} \ill{}
pour $S = \operatorname{Spec} k$ \ill{} \ill{} \ill{} \ill{} $B(k)$, et \ill{}
\ill{} $2$ \ill{} \ill{} \ill{} $\psi(g') - \psi(g)$ ($g', g \in X$) \ill{}}
\note{tout ce paragraphe est encadré et biffé de quatre longues diagonales ; ses dernières lignes sont en outre barrées et surchargées de boucles. On n'en lit que des bribes}

Ceci est raisonnable, grâce au fait que $L$ \ill{} faisceau inv.\ $L_{\beta}$
sur $X$ provient d'un faisceau \struck{\ill{} $L_{\beta}$} $W_{\beta}$ sur $A$
via $X \to A$ \struck{\uncertain{relevé} \ill{} \ill{} $W_{\beta}$}
(\uncertain{défini au moyen} de $g$), savoir le faisceau de Weil $W$ sur
$A \times B$, au cas où $X = A$, $g =$ section unité. Avec les notations
\uncertain{développées} ailleurs, on doit donc \uncertain{vérifier}
\[
{\det} \mathfrak{F}_{\mathcal{L}}(L \otimes L_{\beta}) \simeq {\det} \mathfrak{F}_{\mathcal{L}}(L) \otimes
\underbrace{L'_{-\varphi(\delta)(\beta)}}_{\text{faisceau sur $B$ associé à $-\varphi(\delta)(\beta)$}} .
\]
\marginal{\uncertain{Prouver} \ill{} \uncertain{ces} \uncertain{conditions} \ill{} \ill{}}\note{note oblique dans la marge gauche, en face de « le faisceau de Weil »}
\note{la lettre gothique devant $(L \otimes L_{\beta})$ est lue $\mathfrak{F}$ (transformation de Fourier--Mukai ?), sans certitude ; un premier mot devant « det » est raturé}

\page{19}

Or on a
\struck{$\det$}
\[
\mathfrak{F}_{\mathcal{L}}(L \otimes L_{\beta}) = \mathfrak{F}_{\mathcal{L}}(L) * \varepsilon_{-\beta}
\]
d'où
\[
{\det} \mathfrak{F}_{\mathcal{L}}(L \otimes L_{\beta}) \simeq {\det} \mathfrak{F}_{\mathcal{L}}(L) * \varepsilon_{-\beta}
\]
donc la classe de ${\det}(\mathfrak{F}_{\mathcal{L}}(L \otimes L_{\beta}))$ dans
$\underline{\mathrm{Pic}}_{B/S}(S)$ est égale à \struck{\ill{}} celle de
${\det} \mathfrak{F}_{\mathcal{L}}(L) + \widetilde{\varphi(\delta)}(-\beta)$
(où $\varphi(\delta)$ est l'image de ${\det} \mathfrak{F}_{\mathcal{L}}(L)$ dans
$\mathrm{NS}(B)$), d'où la conclusion annoncée.
\marginal{vérif., \uncertain{avec} \ill{} \uncertain{signes}\ldots}\note{note oblique dans la marge gauche, le dernier mot souligné}

Par suite, on peut interpréter $\ell_g$ comme un hom.\
$\underline{\mathrm{Pic}}^{\delta}_{X/S} \overset{B}{\times} (A, \widetilde{\varphi(\delta)})$
dans $\underline{\mathrm{Pic}}^{\varphi(\delta)}_{B/S}$, ou encore comme
\struck{\ill{}} section de
$\underline{\mathrm{Pic}}^{\delta}_{X/S} \overset{B}{\times} (A, \widetilde{\varphi(\delta)})
\overset{A}{\times} \underline{\mathrm{Pic}}^{\varphi(\delta)}_{B/S}$.
On trouve ainsi \struck{\ill{}} \add{un} \uncertain{morphisme} variable
\[
X \longrightarrow \underline{\mathrm{Pic}}^{\varphi(\delta)}_{B/S} \overset{A}{\times}
\underline{\mathrm{Pic}}^{\delta}_{X/S} \overset{B}{\times} (A, \widetilde{\varphi(\delta)})
\]
qui n'est autre que \struck{\ill{}} l'hom $\ell_{\delta}$ de d), comme on vérifie
facilement.
\note{dans la première formule, un « $-$ » semble précéder le second facteur ($\widetilde{\varphi(\delta)}(-\beta)$) ; la lecture de $\mathfrak{F}_{\mathcal{L}}$ est celle de la p.~18}

\emph{Remarque}\quad a) \uncertain{On} n'a pas pu \ill{}, dans cette
interprétation, à faire intervenir les faisceaux sur
$\underline{\mathrm{Pic}}^{\delta}_{X/S}$, \uncertain{seulement} sur $B$.
Notons que si $\underline{\mathrm{Pic}}^{0}_{X/S}$ \uncertain{admet} pas une
section, \uncertain{on} peut \uncertain{toujours} donner une \uncertain{variante}
aux constructions précédentes, en supposant \emph{donnée} une section \ill{}

\page{20}

$B$ sur $S$ et un hom.\ de groupes
\[
B \xrightarrow{\ i\ } \underline{\mathrm{Pic}}_{X/S} ,
\]
ce qui définit d'ailleurs un torseur $Q$ sous
$A = \underline{\mathrm{Pic}}^{0}_{B/S}$ et un hom canonique
\[
X \xrightarrow{\ \uncertain{\tau}\ } Q
\]
\note{l'étiquette de la flèche est une croix qu'on lit $\tau$ sans certitude ; elle reparaît plus bas}
[C'est par une \uncertain{suite} sur la \uncertain{composante} \ill{} des $i$
et des \uncertain{$\tau$} que devrait \uncertain{commencer} une \uncertain{étude}
\uncertain{de la} petite \uncertain{topologie} de Albanese \ldots].
Tout ce qui a été dit ci-dessus \uncertain{alors}, on trouve
\[
X \times \underline{\mathrm{Pic}}_{X/S} \longrightarrow \underline{\mathrm{Pic}}_{B/S}
\]
définissant
\[
X \times \underline{\mathrm{Pic}}^{\delta}_{X/S} \longrightarrow \underline{\mathrm{Pic}}^{\varphi(\delta)}_{B/S}
\]
\marginal{N.B.\ \ill{} \ill{} avec la notation $\varphi(\delta)$}\note{à droite, souligné « N.B. »}
qui donne en fait, \add{compatible avec $i$}, si \ill{} un \struck{B}-torseur
$P$ et un hom $P \to \underline{\mathrm{Pic}}^{\delta}_{X/S}$,
\[
X \longrightarrow \underline{\mathrm{Pic}}^{\varphi(\delta)}_{B/S} \overset{A}{\times}
\bigl(P \overset{B}{\times} (A, \widetilde{\varphi(\delta)})\bigr)
\]
se factorisant par
\[
Q^{-\chi(\delta)} \longrightarrow \underline{\mathrm{Pic}}^{\varphi(\delta)}_{B/S} \overset{A}{\times}
\bigl(P \overset{B}{\times} (A, \widetilde{\varphi(\delta)})\bigr)
\]
i.e.\ \uncertain{on a} une section de
\[
Q^{\chi(\delta)} \overset{A}{\times} \underline{\mathrm{Pic}}^{\varphi(\delta)}_{B/S}
\overset{A}{\times} \bigl(P \overset{B}{\times} (A, \widetilde{\varphi(\delta)})\bigr)
\]
\note{dans ces trois formules, le premier facteur du produit contracté sous $B$ est écrit $P$ par-dessus un $\underline{\mathrm{Pic}}^{\delta}_{X/S}$ couvert de hachures ; l'exposant de $\underline{\mathrm{Pic}}$ garde un $\delta$ visible}

\struck{b) \ill{} plus \ill{}, dans le \ill{} \ill{} et généralisation en \ill{},
lorsque \ill{} \ill{} \ill{} de sections \ill{} \ill{}, Albanese de \ill{} \ill{}
$X$ \ill{} de sections \ill{} $X \to A$. \ill{}, $X$ \ill{}}
\note{le paragraphe b) de la Remarque, en bas de page, est encadré et biffé de grandes boucles croisées ; la suite est en page 21 (second lot)}

\end{document}
